% Claim proof file for row claim_3_27 -- "LabelRoman".
% Auto-generated stub by scaffold/solve_chapter.py at row start. A manager
% agent dedicated to writing the TeX proof fills in the proof body below.
%
% This file is a sibling of `main.tex` inside `<subsection>/tex/`. The
% `subfiles` package lets it render both standalone and as part of
% `main.tex` via `\subfile{...}`.
%
% The statement is restated above the proof so the file is self-contained
% when rendered on its own. The proof must:
%   - Stay close to the lecture notes' proof if one exists (always search
%     the LN first for a `\begin{proof}` block immediately following the
%     claim; copy / lightly edit when present).
%   - Otherwise use the LN paradigm and cite prior claims by their `ref`.
%   - Be self-contained enough that a mathematician can follow it without
%     re-reading the LN.
%   - Have no `sorry`, `omitted`, or "obvious" shortcuts.

\documentclass[main]{subfiles}

\begin{document}
% ref: claim_3_27 (proof, with statement restated for readability)
% title: LabelRoman
\def\rowref{claim\_3\_27}
\phantomsection\label{claim_3_27}

% --- statement (negation of the positive claim claim_3_27) ------------------
\begin{Lem}[LabelRoman, disproof]\label{lem:replace_walk_disproof}
    The positive lemma stated in the canonical statement file
    \texttt{claim\_3\_27\_statement\_LabelRoman.tex} (the LN's
    \texttt{lem:replace\_walk}) is \emph{false} under the canonical Lean
    encoding of walks, colliders and $\sigma$-openness fixed in
    \texttt{Section3\_1/Walks.lean} (the inductive \texttt{WalkStep} with
    constructors \texttt{forwardE} / \texttt{backwardE} / \texttt{bidir}),
    \texttt{Section3\_3/CollidersAndNon.lean} (the predicates
    \texttt{WalkStep.IsInto} and \texttt{Walk.IsCollider}),
    \texttt{Section3\_3/BlockableAndUnblockable.lean} (the
    \texttt{IsBlockableNonCollider} / \texttt{IsUnblockableNonCollider}
    refinement), and \texttt{Section3\_3/SigmaBlockedWalks.lean}
    (\texttt{Walk.IsSigmaOpenGiven} / \texttt{Walk.IsSigmaBlockedGiven}).
    Precisely, the following counter-existential holds.

    There exist a CDMG $G = (J, V, E, L)$ (def \ref{def-cdmg}; recall that
    def \ref{def-cdmg} permits directed self-loops $(v, v) \in E$ — the
    constraint $E \subseteq (J \cup V) \times V$ has no $v \ne v$
    exclusion — see the chapter-wide registered subtlety
    \texttt{self\_loop\_makes\_tuh\_and\_hut\_simultaneously\_true} in
    \texttt{leanification/working\_subtlety\_register.json}), a subset
    $C \subseteq J \cup V$, a non-negative integer $n \ge 0$, a walk
    \[
        \pi \;=\; (v_0, a_0, v_1, a_1, \dots, v_{n-1}, a_{n-1}, v_n)
    \]
    in $G$ in the sense of def \ref{def:walks}, item~i., that is
    $C$-$\sigma$-open in the sense of def \ref{def:sigma_blocking}, and
    positions $i, j \in \{0, 1, \dots, n\}$ with $i < j$ and $v_i \in
    \Sc^G(v_j)$ (def \ref{def_3_5}), such that the boundary configuration
    of $\pi$ at position $j$ falls into \textbf{case~(i)} of the LN's
    replacement procedure (i.e., $j = n$ or $a_j = (v_j, v_{j+1}) \in E$),
    and such that for the unique shortest directed walk
    \[
        \sigma_{ij} \;=\; (w_0 = v_i, \, b_0, \, w_1, \, b_1, \, \dots, \,
        b_{m-1}, \, w_m = v_j)
    \]
    from $v_i$ to $v_j$ in $G$ (def \ref{def:walks}, item~ii.) prescribed
    by case~(i), the modified walk
    \[
        \pi' \;=\; \bigl(v_0, \, a_0, \, \dots, \, a_{i-1}, \, v_i = w_0, \,
        b_0, \, \dots, \, b_{m-1}, \, w_m = v_j, \, a_j, \, \dots, \,
        v_n\bigr)
    \]
    obtained from $\pi$ by replacing the subwalk between positions $i$ and
    $j$ by $\sigma_{ij}$ is $C$-$\sigma$-\emph{blocked} in the sense of def
    \ref{def:sigma_blocking}.

    Consequently, conclusion~(2) of the canonical positive statement — ``the
    modified walk $\pi'$ is $C$-$\sigma$-open'' — fails on this instance, and
    \emph{a fortiori} the conjunction of conclusions~(1) and~(2) of the
    canonical positive statement is false. The negation of the positive
    claim is thereby witnessed.
\end{Lem}

% --- proof ------------------------------------------------------------------
\begin{proof}
We exhibit an explicit counter-example. The witness was independently
verified end-to-end by a downstream \texttt{verify\_tex\_proof} worker
session (session id \texttt{0b28bf97-9527-4042-b605-3bca225ff00f}); the
proof below re-derives every required predicate value directly from the
canonical Lean encoding cited in the statement so that the failure is
manifestly traceable to the encoding-level facts.

\medskip\noindent\textbf{Step 1 — the CDMG $G$.}
Let $G = (J, V, E, L)$ with
\begin{align*}
    J &\;=\; \emptyset, \\
    V &\;=\; \{a, b\}, \\
    E &\;=\; \{\, (a, b),\; (b, a),\; (b, b) \,\} \;\subseteq\; (J \cup V) \times V, \\
    L &\;=\; \emptyset \;\subseteq\; \{\, \{x, y\} : x, y \in J \cup V \,\}.
\end{align*}
This is a CDMG in the sense of def \ref{def-cdmg}: $J$, $V$ are disjoint
finite sets, $E$ consists of ordered pairs whose head lies in $V$ (the
constraint $E \subseteq (J \cup V) \times V$), and $L$ is a set of
unordered pairs in $J \cup V$. The directed self-loop $(b, b)$ is admitted
because def \ref{def-cdmg}'s definition of $E$ has no $v \ne v$
exclusion — it is exactly this admissibility (chapter-wide registered
subtlety \texttt{self\_loop\_makes\_tuh\_and\_hut\_simultaneously\_true})
that the counter-example exploits.

\medskip\noindent\textbf{Step 2 — strongly connected components and the subset $C$.}
The directed walk $(a, (a, b), b)$ in $G$ (length~$1$, using $(a, b) \in E$
forward) witnesses $a \in \Anc^G(b)$ per def \ref{def_3_5}, item~iv.
Symmetrically, the directed walk $(b, (b, a), a)$ in $G$ (length~$1$,
using $(b, a) \in E$ forward) witnesses $a \in \Desc^G(b)$, i.e.\
$b \in \Anc^G(a)$, equivalently $a \in \Desc^G(b)$. Hence
\[
    a \;\in\; \Anc^G(b) \cap \Desc^G(b) \;=\; \Sc^G(b).
\]
Combined with the unconditional self-membership $b \in \Sc^G(b)$ (def
\ref{def_3_5}, via the trivial directed walk $(b)$ of length $0$ at $b$,
per def \ref{def:walks}, item~ii.\ with empty edge list), and the fact
that $V = \{a, b\}$ exhausts the candidates, we have
\[
    \Sc^G(b) \;=\; \{a, b\}.
\]
The self-loop $(b, b) \in E$ is consistent with this computation but is
not needed for it — the cycle $a \tuh b \tuh a$ already establishes the
SCC.

Let $C := \emptyset \subseteq J \cup V$. Then
\[
    \Anc^G(C) \;=\; \bigcup_{c \in C} \Anc^G(c) \;=\; \emptyset
\]
(def \ref{def_3_5}, item~iv., applied to the empty index set; in the Lean
formalisation, \texttt{AncSet} of \texttt{Section3\_1/Family\-Relation\-ships.lean}
line~204 specialises to $\bigcup_{v \in \emptyset} G.\mathtt{Anc}\, v
= \emptyset$).

\medskip\noindent\textbf{Step 3 — the walk $\pi$ in $G$.}
Let $n = 2$ and define $\pi = (v_0, a_0, v_1, a_1, v_2)$ with
\begin{itemize}
    \item vertex sequence $(v_0, v_1, v_2) = (a, b, b)$,
    \item walk-edge $a_0 := (b, a) \in E$, used as the \emph{second}
        disjunct of the walk constraint of def \ref{def:walks}, item~i.,
        at index $0$: $a_0 = (v_1, v_0) = (b, a) \in E$. Visually
        $v_0 \hut v_1$ — the head of the directed edge $(b, a)$ is at the
        LEFT vertex $v_0 = a$, its tail is at the RIGHT vertex
        $v_1 = b$. In the Lean encoding (\texttt{Section3\_1/Walks.lean}
        lines~278--282, inductive \texttt{WalkStep}) this is
        \texttt{backwardE}\,$h_{ba}$ of type \texttt{WalkStep G a b}
        carrying the membership witness $h_{ba} : (b, a) \in G.\mathtt{E}$;
    \item walk-edge $a_1 := (b, b) \in E$, used as the \emph{first}
        disjunct of the walk constraint at index $1$: $a_1 = (v_1, v_2) =
        (b, b) \in E$. This is the directed self-loop. In the Lean
        encoding this is \texttt{forwardE}\,$h_{bb}$ of type
        \texttt{WalkStep G b b} carrying the membership witness
        $h_{bb} : (b, b) \in G.\mathtt{E}$.
\end{itemize}
The walk constraint of def \ref{def:walks}, item~i., is satisfied at both
indices, so $\pi$ is a walk in $G$ in the sense of def \ref{def:walks},
item~i. Its length is $n = 2$ and
$\pi.\mathtt{vertices} = [a, b, b]$ (computed by the Lean recursion of
\texttt{Section3\_1/Walks.lean} lines~594--597).

\medskip\noindent\textbf{Step 4 — $\pi$ is $C$-$\sigma$-open.}
We check both clauses of def \ref{def:sigma_blocking} (Lean encoding:
\texttt{Walk.IsSigmaOpenGiven} at \texttt{Section3\_3/SigmaBlockedWalks.lean}
lines~366--372) at every position $k \in \{0, 1, 2\}$ on $\pi$.

\emph{Position $k = 0$ ($v_0 = a$, left endnode).}
\texttt{Walk.IsCollider} (\texttt{Section3\_3/CollidersAndNon.lean}
lines~628--635) fires \texttt{False} at $k = 0$ via the
\texttt{.cons \_ \_ (.cons \_ \_ \_), 0} branch — end-positions are not
colliders. So the first conjunct of \texttt{IsSigmaOpenGiven} is vacuous
at $k = 0$. \texttt{Walk.IsBlockableNonCollider} at $k = 0$
(\texttt{Section3\_3/BlockableAndUnblockable.lean} lines~608--613) reduces
to \texttt{IsNonCollider 0 $\land$ (0 = 0 $\lor$ \dots) = True}, so the
second conjunct of \texttt{IsSigmaOpenGiven} at $k = 0$ requires
$v_0 = a \notin C = \emptyset$, which holds.

\emph{Position $k = 2$ ($v_2 = b$, right endnode).}
Symmetric to position~$0$: the recursion of \texttt{IsCollider} descends
to a \texttt{.nil}-tail branch and returns \texttt{False}, and
\texttt{IsBlockableNonCollider} fires the $k = \mathtt{p.length}$ disjunct
(here $\mathtt{p.length} = 2 = k$). The second conjunct of
\texttt{IsSigmaOpenGiven} at $k = 2$ requires $v_2 = b \notin C =
\emptyset$, which holds.

\emph{Position $k = 1$ ($v_1 = b$, interior — the load-bearing case).}
We compute \texttt{Walk.IsCollider $\pi$ 1} by unfolding the cons-cell
pattern:
\[
    \pi \;=\; \mathtt{cons}\;b\;(\mathtt{backwardE}\,h_{ba})\;
              \bigl(\mathtt{cons}\;b\;(\mathtt{forwardE}\,h_{bb})\;
              (\mathtt{nil}\;b\;h_b)\bigr).
\]
The pattern \texttt{.cons vk s$_0$ (.cons \_ s$_1$ \_), 1} matches with
$\mathtt{vk} = b$, $s_0 = \mathtt{backwardE}\,h_{ba}$ of type
\texttt{WalkStep G a b}, and $s_1 = \mathtt{forwardE}\,h_{bb}$ of type
\texttt{WalkStep G b b}. The body evaluates to $s_0.\mathtt{IsInto}\,b
\;\land\; s_1.\mathtt{IsInto}\,b$ where \texttt{WalkStep.IsInto}
(\texttt{Section3\_3/CollidersAndNon.lean} lines~338--343) reads:
\begin{itemize}
    \item For $s_0 = \mathtt{backwardE}\,h_{ba}$ with implicit type indices
        $(u, v) = (a, b)$: \texttt{IsInto $w$} unfolds to
        ``$w = u \lor (s(u, v) \in G.L \land (w = u \lor w = v))$''.
        Substitute $w := b$, $u := a$, $v := b$: $b = a$ is \textbf{false};
        $s(a, b) \in G.L$ requires $\{a, b\} \in L = \emptyset$, also
        \textbf{false}. Hence $s_0.\mathtt{IsInto}\,b =
        \mathtt{False}$.
    \item For $s_1 = \mathtt{forwardE}\,h_{bb}$ with implicit type indices
        $(u, v) = (b, b)$: \texttt{IsInto $w$} unfolds to
        ``$w = v \lor (s(u, v) \in G.L \land (w = u \lor w = v))$''.
        Substitute $w := b$, $v := b$: $b = b$ is \textbf{true}. Hence
        $s_1.\mathtt{IsInto}\,b = \mathtt{True}$.
\end{itemize}
Therefore $\mathtt{Walk.IsCollider}\;\pi\;1 = \mathtt{False} \land
\mathtt{True} = \mathtt{False}$: position~$1$ is a \emph{non-collider}
on~$\pi$, with arrowhead count $\mathrm{ah}_\pi(1) = 0 + 1 = 1$. The
first conjunct of \texttt{IsSigmaOpenGiven} at $k = 1$ is therefore
vacuous.

For the second conjunct of \texttt{IsSigmaOpenGiven} at $k = 1$, the
requirement is ``$v_1 \notin C$ whenever $\pi$ has
\texttt{IsBlockableNonCollider} at~$1$''. Either way (whether
position~$1$ is blockable or unblockable), $v_1 = b \notin C = \emptyset$
holds vacuously, so the clause is satisfied. (For completeness:
$\mathtt{HasBlockingLeftSlot}\;\pi\;1$ unfolds to
``$a \notin G.\mathtt{Sc}\,b$'' — \textbf{false}, since $a \in \Sc^G(b)
= \{a, b\}$ by Step~2. $\mathtt{HasBlockingRightSlot}\;\pi\;1$ recurses
to the tail $\mathtt{cons}\;b\;(\mathtt{forwardE}\,h_{bb})\;\mathtt{nil}$
at $k = 0$ and unfolds to ``$b \notin G.\mathtt{Sc}\,b$'' —
\textbf{false}, by $\Sc$-self-reflexivity. So $v_1$ is in fact an
\emph{unblockable} non-collider on $\pi$;
\texttt{IsBlockableNonCollider} fires \texttt{False} and the second
conjunct of \texttt{IsSigmaOpenGiven} is vacuous at $k = 1$.)

All three positions satisfy both clauses of
\texttt{Walk.IsSigmaOpenGiven}, so $\pi$ is $C$-$\sigma$-open. $\;\;\checkmark$

\medskip\noindent\textbf{Step 5 — the hypothesis on positions $(i, j)$.}
Take $i := 0$ and $j := 1$. Then
\[
    0 \;=\; i \;<\; j \;=\; 1 \;\le\; 2 \;=\; n,
\]
and $v_i = v_0 = a \in \Sc^G(b) = \Sc^G(v_j)$ by Step~2. The hypothesis
of the positive lemma on $(i, j)$ is satisfied. Note that $v_i = a \ne
b = v_j$, so this is \emph{not} the degenerate $v_i = v_j$ branch
addressed by the canonical statement's ``Addition to the LN'' paragraph.

\medskip\noindent\textbf{Step 6 — case~(i) of the LN procedure fires.}
The case-(i) trigger of the canonical positive statement reads
``$j = n \;\lor\; a_j = (v_j, v_{j+1}) \in E$''.
Here $j = 1 \ne 2 = n$, so the first disjunct is false.
For the second disjunct: $a_j = a_1$ is the self-loop step described in
Step~3, with $a_1 = (v_1, v_2) = (b, b) \in E$ (the first-disjunct
encoding of the walk constraint at index~$1$). So the second disjunct
holds and case~(i) is the prescribed branch.

\medskip\noindent\textbf{Step 7 — the unique shortest directed walk
$\sigma_{ij}$ from $v_i = a$ to $v_j = b$ in $G$.}
We enumerate directed walks (def \ref{def:walks}, item~ii.) from $a$ to
$b$ in $G$ by length:
\begin{itemize}
    \item Length $0$: the trivial directed walk $(u_0)$ at $u_0 = u = w$
        requires $a = u = w = b$, which fails since $a \ne b$. So no
        length-$0$ directed walk from $a$ to $b$ exists in $G$.
    \item Length $1$: a directed walk $(a, b_0, b)$ with $b_0 = (a, b)
        \in E$. The unique directed edge in $E$ with tail $a$ and head $b$
        is $(a, b)$ itself, so up to the data of its single step there is
        \emph{exactly one} length-$1$ directed walk from $a$ to $b$ in
        $G$, namely
        \[
            \sigma \;:=\; (a, (a, b), b).
        \]
    \item Length $\ge 2$: any such directed walk has length strictly
        greater than $1$, hence is not shortest.
\end{itemize}
Therefore the shortest directed walk from $v_i = a$ to $v_j = b$ in $G$
has length $m = 1$ and is unique:
\[
    \sigma_{ij} \;=\; (w_0 = a,\; b_0 = (a, b),\; w_1 = b).
\]
In the Lean encoding this is the walk $\mathtt{cons}\;b\;(\mathtt{forwardE}\,
h_{ab})\;(\mathtt{nil}\;b\;h_b)$ where $h_{ab} : (a, b) \in G.\mathtt{E}$
witnesses the directed edge, of type \texttt{Walk G a b}.

\medskip\noindent\textbf{Step 8 — the modified walk $\pi'$.}
Per the canonical statement's prescription,
$\pi' = (v_0, a_0, \dots, a_{i-1}, v_i = w_0, b_0, \dots, b_{m-1}, w_m =
v_j, a_j, \dots, v_n)$. With $i = 0$ (the prefix $(v_0, a_0, \dots,
a_{i-1}, v_i)$ collapses to the single-vertex $(v_0) = (a)$), $m = 1$,
and $j = 1$ (the suffix $(v_j, a_j, v_{j+1}, \dots, v_n)$ is
$(v_1, a_1, v_2) = (b, (b, b), b)$), we obtain
\[
    \pi' \;=\; \bigl(v'_0 = a,\;\; a'_0 = (a, b) \in E,\;\;
                    v'_1 = b,\;\; a'_1 = (b, b) \in E,\;\;
                    v'_2 = b\bigr).
\]
Both walk-steps of $\pi'$ are first-disjunct encodings of def
\ref{def:walks}, item~i.:
\begin{itemize}
    \item $a'_0 = (v'_0, v'_1) = (a, b) \in E$ (the new step from
        $\sigma_{ij}$);
    \item $a'_1 = (v'_1, v'_2) = (b, b) \in E$ (the self-loop step
        inherited verbatim from $\pi$'s slot~$j = 1$ — the replacement
        only edits the subwalk between positions $i = 0$ and $j = 1$,
        which contains only $a_0$).
\end{itemize}
In the Lean encoding,
\[
    \pi' \;=\; \mathtt{cons}\;b\;(\mathtt{forwardE}\,h_{ab})\;
              \bigl(\mathtt{cons}\;b\;(\mathtt{forwardE}\,h_{bb})\;
              (\mathtt{nil}\;b\;h_b)\bigr).
\]
Both $\pi$ and $\pi'$ have the same vertex sequence $[a, b, b]$ and the
same length $2$; they differ \emph{only} in the encoding of the first
walk-step (\texttt{backwardE} on $\pi$, \texttt{forwardE} on $\pi'$).

\medskip\noindent\textbf{Step 9 — $\pi'$ is $C$-$\sigma$-blocked at
position~$1$.}
We recompute \texttt{Walk.IsCollider} at position~$1$ on $\pi'$:
\[
    \pi' \;=\; \mathtt{cons}\;b\;(\mathtt{forwardE}\,h_{ab})\;
              \bigl(\mathtt{cons}\;b\;(\mathtt{forwardE}\,h_{bb})\;
              (\mathtt{nil}\;b\;h_b)\bigr).
\]
The pattern \texttt{.cons vk s$'_0$ (.cons \_ s$'_1$ \_), 1} matches with
$\mathtt{vk} = b$, $s'_0 = \mathtt{forwardE}\,h_{ab}$ of type
\texttt{WalkStep G a b}, and $s'_1 = \mathtt{forwardE}\,h_{bb}$ of type
\texttt{WalkStep G b b}. The body evaluates to $s'_0.\mathtt{IsInto}\,b
\;\land\; s'_1.\mathtt{IsInto}\,b$ where (again by
\texttt{Section3\_3/CollidersAndNon.lean} lines~338--343):
\begin{itemize}
    \item For $s'_0 = \mathtt{forwardE}\,h_{ab}$ with type indices
        $(u, v) = (a, b)$: \texttt{IsInto $b$} unfolds to ``$b = v \lor
        \dots$''; substitute $v := b$, giving $b = b = \mathtt{True}$.
        Hence $s'_0.\mathtt{IsInto}\,b = \mathtt{True}$.
    \item For $s'_1 = \mathtt{forwardE}\,h_{bb}$ with type indices
        $(u, v) = (b, b)$: \texttt{IsInto $b$} unfolds as in Step~4 to
        $b = b = \mathtt{True}$. Hence $s'_1.\mathtt{IsInto}\,b =
        \mathtt{True}$.
\end{itemize}
Therefore $\mathtt{Walk.IsCollider}\;\pi'\;1 = \mathtt{True} \land
\mathtt{True} = \mathtt{True}$: position~$1$ is a \emph{collider} on
$\pi'$, with arrowhead count $\mathrm{ah}_{\pi'}(1) = 1 + 1 = 2$.

By def \ref{def:sigma_blocking} (Lean encoding:
\texttt{Walk.IsSigma\-Open\-Given} at
\texttt{Section3\_3/SigmaBlockedWalks.lean} lines~366--372), the first
conjunct of $C$-$\sigma$-openness at position~$1$ requires
\[
    v'_1 \;\in\; \Anc^G(C) \;=\; \Anc^G(\emptyset) \;=\; \emptyset.
\]
But $v'_1 = b \notin \emptyset$. So $\pi'$ \emph{violates} the
$\sigma$-open collider clause at position~$1$, and dually (Lean encoding:
\texttt{Walk.IsSigmaBlockedGiven} at
\texttt{Section3\_3/SigmaBlockedWalks.lean} lines~512--518), the
existential witness
\[
    \bigl\langle k = 1,\;\; v'_1 = b,\;\;
        \pi'.\mathtt{vertices}[1]? = \mathtt{some}\,b,\;\;
        \pi'.\mathtt{IsCollider}\,1,\;\;
        b \notin \mathtt{G.AncSet}\,\emptyset \bigr\rangle
\]
closes the first existential disjunct of
$\pi'.\mathtt{IsSigmaBlockedGiven}\;\emptyset\;h_C$ (where
$h_C : \emptyset \subseteq \uparrow G.\mathtt{J} \cup \uparrow
G.\mathtt{V}$ is the vacuous subset inclusion). Hence $\pi'$ is
$C$-$\sigma$-blocked. $\;\;\checkmark$

\medskip\noindent\textbf{Step 10 — synthesis.}
We have constructed:
\begin{itemize}
    \item a CDMG $G$ with $J = \emptyset$, $V = \{a, b\}$, $E = \{(a, b),
        (b, a), (b, b)\}$, $L = \emptyset$ (Step~1);
    \item a subset $C = \emptyset \subseteq J \cup V$ (Step~2);
    \item a walk $\pi$ in $G$ of length $n = 2$ with vertex sequence
        $(a, b, b)$ (Step~3) that is $C$-$\sigma$-open (Step~4);
    \item positions $i = 0$, $j = 1$ with $i < j$ and $v_i \in \Sc^G(v_j)$
        (Step~5);
    \item the case-(i) trigger of the LN procedure fires at position $j$
        (Step~6);
    \item the unique shortest directed walk $\sigma_{ij} = (a, (a, b), b)$
        from $v_i$ to $v_j$ in $G$ (Step~7);
    \item the modified walk $\pi' = (a, (a, b), b, (b, b), b)$ (Step~8);
\end{itemize}
and we have shown that $\pi'$ is $C$-$\sigma$-\emph{blocked} (Step~9).
This is a counter-example to conclusion~(2) of the canonical positive
statement (``$\pi'$ is $C$-$\sigma$-open''), and therefore witnesses the
negation of the positive lemma as stated above.

\medskip\noindent\textbf{Remark on the mechanism (for downstream
remediation).}
The discrepancy between the LN's pictorial reading (which treats
position~$1$ on $\pi'$ as a non-collider via the configuration
$v'_0 \tuh v'_1 \tuh v'_2$ where the tail of the self-loop at $v'_1$
``cancels'' its head, by viewing each step's boundary mark at $v'_1$
only) and the canonical Lean encoding's node-equality reading of
\texttt{IsInto} (which fires on \emph{any} type-index match, so the
self-loop's head at $v_2 = b$ also counts as an arrowhead at $v_1 = b$
since $v_1 = v_2 = b$) is the load-bearing source of the bug. The
chapter-wide registered subtlety
\texttt{self\_loop\_makes\_tuh\_and\_hut\_simultaneously\_true} (recorded
against def \ref{def:collider_noncollider} in
\texttt{leanification/working\_subtlety\_register.json}) documents the
encoding decision; the consumer-side propagation into the present bug
is recorded under the id \texttt{claim\_3\_27\_case\_i\_fails\_for\_self\_loops}
in the same register, together with a verifier-recommended remediation
that tightens case~(i)'s trigger to ``$j = n \;\lor\; (a_j = (v_j, v_{j+1})
\in E \;\land\; v_{j+1} \ne v_j)$'', re-routing self-loops at step $j$
into case~(ii) where they are absorbed without introducing a new collider.
The present disproof file establishes the negation under the canonical
encoding without that tightening; whether the codebase adopts the
remediation via an \texttt{addition\_to\_the\_LN} on the row, or accepts
the disprove verdict and proves only the negation, is an orchestrator-
level decision recorded in \texttt{leanification/working\_subtlety\_register.json}.
\end{proof}

\end{document}
